                   IMC 2015, Blagoevgrad, Bulgaria

                               Day 1, July 29, 2015




Problem 1.     For any integer n ≥ 2 and two n × n matrices with real entries A, B that
satisfy the equation
                                A−1 + B −1 = (A + B)−1
prove that det(A) = det(B).
   Does the same conclusion follow for matrices with complex entries?
                  (Proposed by Zbigniew Skoczylas, Wrocªaw University of Technology)
Solution.   Multiplying the equation by (A + B) we get
                    I = (A + B)(A + B)−1 = (A + B)(A−1 + B −1 ) =
                = AA−1 + AB −1 + BA−1 + BB −1 = I + AB −1 + BA−1 + I
                                AB −1 + BA−1 + I = 0.

Let X = AB −1 ; then A = XB and BA−1 = X −1 , so we have X +X −1 +I = 0; multiplying
by (X − I)X ,
            0 = (X − I)X · (X + X −1 + I) = (X − I) · (X 2 + X + I) = X 3 − I.

Hence,
                                         X3 = I
                             (det X)3 = det(X 3 ) = det I = 1
                                       det X = 1
                        det A = det(XB) = det X · det B = det B.
                                                                         √
    In case of complex matrices the statement is false. Let ω = 21 (−1 + i 3). Obviously
ω∈ / R and ω 3 = 1, so 0 = 1 + ω + ω 2 = 1 + ω + ω .
    Let A = I and let B be a diagonal matrix with all entries along the diagonal equal to
either ω or ω = ω 2 such a way that det(B) 6= 1 (if n is not divisible by 3 then one may
set B = ωI ). Then A−1 = I , B −1 = B . Obviously I + B + B = 0 and
                    (A + B)−1 = (−B)−1 = −B = I + B = A−1 + B −1 .

By the choice of A and B , det A = 1 6= det B .
Problem 2.       For a positive integer n, let f (n) be the number obtained by writing n
in binary and replacing every 0 with 1 and vice versa. For example, n = 23 is 10111 in
binary, so f (n) is 1000 in binary, therefore f (23) = 8. Prove that
                                              n
                                              X          n2
                                                  f (k) ≤ .
                                              k=1
                                                         4

When does equality hold?
                                         (Proposed by Stephan Wagner, Stellenbosch University)
Solution.    If r and k are positive integers with 2r−1 ≤ k < 2r then k has r binary digits,
so k + f (k) = 11  . . . 1} = 2r − 1.
                           (2)
                | {z
                   r
   Assume that 2s−1 − 1 ≤ n ≤ 2s − 1. Then
                                        n           n
                             n(n + 1) X            X
                                     +     f (k) =     (k + f (k)) =
                                2      k=1         k=1
                           s−1
                           X         X                        X
                       =                    (k + f (k)) +           (k + f (k)) =
                           r=1 2r−1 ≤k<2r                   2s−1 ≤k≤n
                           s−1
                           X
                       =         2r−1 · (2r − 1) + (n − 2s−1 + 1) · (2s − 1) =
                           r=1
                           s−1
                           X                s−1
                                            X
                                 2r−1
                       =         2      −         2r−1 + (n − 2s−1 + 1)(2s − 1) =
                       r=1       r=1
               = 3 (4 − 1) − (2 − 1) + (2s − 1)n − 22s−1 + 3 · 2s−1 − 1 =
                 2 s−1         s−1

                              = (2s − 1)n − 31 4s + 2s − 32

and therefore
                   n
              n2 X            n2
                                                                    
                                     s        1 s    s  2   n(n + 1)
                −     f (k) =    − (2 − 1)n − 3 4 + 2 − 3 −            =
              4   k=1
                              4                                2
                              = 34 n2 − (2s − 23 )n + 13 4s − 2s + 23 =
                                        2s+1 − 2             2s+1 − 4
                                                                   
                               3
                             =      n−                 n−               .
                               4            3                    3

Notice that the dierence of the last two factors is less than 1, and one of them must
be an integer: 2 3 −2 is integer if s is even, and 2 3 −4 is integer if s is odd. Therefore,
                 s+1                                s+1


either one of them is 0, resulting a zero product, or both factors have the same sign, so
the product is strictly positive. This solves the problem and shows that equality occurs
         2s+1 − 2                    2s+1 − 4
if n =            (s is even) or n =          (s is odd).
            3                           3
Problem 3.     Let F (0) = 0, F (1) = 32 , and F (n) = 25 F (n − 1) − F (n − 2) for n ≥ 2.
                                  ∞
                                          1
     Determine whether or not                  is a rational number.
                                  X

                                   n=0
                                       F (2n )
                    (Proposed by Gerhard Woeginger, Eindhoven University of Technology)
Solution 1.    The characteristic equation of our linear recurrence is x2 − 52 x + 1 = 0, with
roots x1 = 2 and x2 = 21 . So F (n) = a · 2n + b · ( 21 )n with some constants a, b. By F (0) = 0
and F (1) = 32 , these constants satisfy a + b = 0 and 2a + 2b = 32 . So a = 1 and b = −1,
and therefore
                                             F (n) = 2n − 2−n .
     Observe that
                             n
               1         22        1      1        1      1
                 n
                   =   2n 2
                              = 2n   − 2n 2   = 2n   − 2n+1    ,
            F (2 )   (2 ) − 1  2 − 1 (2 ) − 1  2 −1 2       −1
so                ∞            ∞                 
                  X      1     X      1      1           1
                           n
                             =     2n   − 2n+1      = 20   = 1.
                  n=0
                      F (2 ) n=0 2 − 1 2       −1    2 −1
Hence the sum takes the value 1, which is rational.
Solution 2. As in the rst solution we nd that F (n) = 2 − 2   . Then
                                                         n   −n

              ∞              ∞                ∞             n
              X      1      X        1       X      ( 12 )2
                          =                =
              n=0
                  F (2n )   n=0
                                22n − 2−2n   n=0
                                                 1 − ( 12 )2n+1
                                 ∞               ∞                k       ∞               ∞
                                             n               n+1                        n                  n
                                 X               X                          X               X
                           =           ( 12 )2         ( 21 )2          =         ( 12 )2         ( 12 )2k·2
                                 n=0             k=0                        n=0             k=0
                                 ∞ X
                                   ∞                                ∞
                                                   n
                                 X                                  X
                           =                 ( 12 )2 (2k+1) =             ( 12 )m = 1.
                                 n=0 k=0                            m=1

(Here we used the fact that every positive integer m has a unique representation m =
2n (2k + 1) with non-negative integers n and k .)
    This shows that the series converges to 1.


Problem 4.     Determine whether or not there exist 15 integers m1 , . . . , m15 such that
                                 15
                                 X
                                       mk · arctan(k) = arctan(16).                                            (1)
                                 k=1

                    (Proposed by Gerhard Woeginger, Eindhoven University of Technology)
Solution. We show that such integers m1 , . . . , m15 do not exist.
   Suppose that (1) is satised by some integers m1 , . . . , m15 . Then the argument of the
complex number z1 = 1 + 16i coincides with the argument of the complex number
                  z2 = (1 + i)m1 (1 + 2i)m2 (1 + 3i)m3 · · · · · · (1 + 15i)m15 .

Therefore the ratio R = z2 /z1 is real (and not zero). As Re z1 = 1 and Re z2 is an integer,
R is a nonzero integer.
      By considering the squares of the absolute values of z1 and z2 , we get
                                                                      15
                                                                      Y
                                                      2      2
                                             (1 + 16 )R           =        (1 + k 2 )mk .
                                                                      k=1

Notice that p = 1 + 162 = 257 is a prime (the fourth Fermat prime), which yields an
easy contradiction through p-adic valuations: all prime factors in the right hand side are
strictly below p (as k < 16 implies 1 + k2 < p). On the other hand, in the left hand side
the prime p occurs with an odd exponent.


Problem 5.        Let n ≥ 2, let A1 , A2 , . . . , An+1 be n + 1 points in the n-dimensional
Euclidean space, not lying on the same hyperplane, and let B be a point strictly inside
the convex hull of A1 , A2 , . . . , An+1 . Prove that ∠Ai BAj > 90◦ holds for at least n pairs
(i, j) with 1 ≤ i < j ≤ n + 1.
                                        (Proposed by Géza Kós, Eötvös University, Budapest)
                              −−→
Solution.  Let vi = BAi . The condition ∠Ai BAj > 90◦ is equivalent with vi · vj < 0.
Since B is an interior point of the simplex, there are some weights w1 , . . . , wn+1 > 0 with
n+1
      w i v i = 0.
P
i=1
    Let us build a graph on the vertices 1, . . . , n + 1. Let the vertices i and j be connected
by an edge if vi · vj < 0. We show that this graph is connected. Since every connected
graph on n + 1 vertices has at least n edges, this will prove the problem statement.
    Suppose the contrary that the graph is not connected; then the vertices can be split
in two disjoint nonempty sets, say V and W such that V ∪ W = {1, 2, . . . , n + 1}. Since
there is no edge between the two vertex sets, we have vi · vj ≥ 0 for all i ∈ V and j ∈ W .
    Consider
                                    !2                       !2                      !2
                     X                        X                       X                        XX
            0=             w i vi        =          w i vi        +          wi vi        +2             wi wj (vi · vj ).
                 i∈V ∪W                       i∈V                     i∈W                      i∈V i∈W

Notice that all terms are nonnegative on the right-hand side. Moreover,                                           wi vi 6= 0 and
                                                                                                             P
                                                                                                            i∈V
       wi vi 6= 0, so there are at least two strictly nonzero terms, contradiction.
P
i∈W



Remark 1.              The number n in the statement is sharp; if vn+1 = (1, 1, . . . , 1) and vi =

(0, . . . , 0, −1, 0, . . . , 0) for i = 1, . . . , n then vi · vj < 0 holds only when i = n + 1 or j = n + 1.
 | {z }            | {z }
      i−1            n−i



                                       http://math.stackexchange.com/questions/476640/n
Remark 2. The origin of the problem is here:
-simplex-in-an-intersection-of-n-balls/789390
